\documentclass{article}
\usepackage{amsmath,amssymb,amsthm}
\usepackage[margin=1in]{geometry}

\newcommand{\area}{\operatorname{area}}
\newcommand{\dinv}{\operatorname{dinv}}
\newcommand{\defc}{\operatorname{defc}}
\newcommand{\Cat}{\mathcal C}
\newcommand{\D}{\mathcal D}
\newcommand{\T}{\mathcal T}
\newtheorem{conjecture}{Conjecture}

\title{A Conjectural Formula for the Rational \(q,t\)-Catalan Polynomial}
\author{}
\date{}

\begin{document}
\maketitle

\section{Purpose and status}

Fix coprime integers \(r,s\ge2\).  The rational \(q,t\)-Catalan polynomial
is a generating function over rational Dyck paths from \((0,0)\) to
\((s,r)\).  Hawkes proposed a formula that expresses this polynomial as a
signed sum of symmetric monomial strings indexed by a distinguished subset of
paths~\cite{Hawkes2024}.  The formula is attractive because its right-hand
side is symmetric in \(q\) and \(t\) by construction.

The formula remains conjectural in general.  The accompanying
\texttt{code.py} performs exact finite checks: it enumerates all rational Dyck
paths for a chosen coprime pair, forms both coefficient dictionaries, and
compares them.  Thus a passing run proves the finite identity for that pair,
but does not prove the conjecture uniformly for all \((r,s)\).

\section{Position-coordinate paths}

For \(0\le i<s\), define the ceiling and label of column \(i\) by
\[
  H_i=\left\lfloor\frac{ri}{s}\right\rfloor,
  \qquad
  L_i\equiv ri\pmod s,
  \qquad 0\le L_i<s.
\]
Because \(r\) and \(s\) are coprime, the labels are a permutation of
\(0,1,\ldots,s-1\).

A position-coordinate path is an integer sequence
\[
  Q=(Q_0,\ldots,Q_{s-1}),
  \qquad
  Q_0=0,
  \qquad
  0\le Q_i\le H_i.
\]
Put
\[
  P_i=H_i-Q_i.
\]
The sequence represents an \((r,s)\)-Dyck path when
\[
  P_0\le P_1\le\cdots\le P_{s-1}.
\]
The area statistic is particularly simple in these coordinates:
\[
  \area(Q)=\sum_{i=0}^{s-1}Q_i.
\]
The maximum total degree is
\[
  M=\sum_{i=0}^{s-1}H_i
   =\frac{(r-1)(s-1)}{2}.
\]

For completeness, we briefly describe the deficit statistic used by the
checker.  Set
\[
  q_i=Q_i+\frac{L_i}{s},
  \qquad
  w_i=P_{i+1}-P_i
  \quad(0\le i<s-1).
\]
For \(1\le i<j<s\), define
\[
  \delta_{ij}(Q)
  =
  \sum_{n=1}^{w_i}{\bf 1}(q_j\le q_i-n)
  +
  \sum_{n=1}^{w_{i-1}}{\bf 1}(q_j\ge q_i+n),
\]
where the first sum is absent when \(i=s-1\), a case that does not occur
because \(i<j<s\).  Then
\[
  \defc(Q)=\sum_{1\le i<j<s}\delta_{ij}(Q),
  \qquad
  \dinv(Q)=M-\area(Q)-\defc(Q).
\]
The code evaluates an equivalent integer formula, avoiding floating-point
arithmetic.

Let \(\D_{r/s}\) denote the set of these paths.  The rational
\(q,t\)-Catalan polynomial is
\[
  \Cat_{r/s}(q,t)
  =
  \sum_{Q\in\D_{r/s}}q^{\area(Q)}t^{\dinv(Q)}.
\]

\section{The distinguished paths}

There is a unique column \(c\) with
\[
  L_c=1.
\]
Since
\[
  \frac{rc}{s}=H_c+\frac1s,
\]
the lattice point \((c,H_c)\) is the lattice point immediately below the
bounding diagonal with the smallest possible positive gap.  This is the
``closest point'' used in the paper.

Define
\[
  \T_{r/s}
  =
  \{Q\in\D_{r/s}:Q_c=0\}.
\]
The condition \(Q_c=0\) says exactly that the path passes through
\((c,H_c)\).  These are the maximal or distinguished paths that index the
conjectural formula.

\section{The signed string formula}

For nonnegative integers \(a,b\), define
\[
  \operatorname{Str}(a,b)=
  \begin{cases}
    \displaystyle\sum_{j=a}^{b}q^j t^{a+b-j},&a\le b,\\[3mm]
    \displaystyle-\sum_{j=b+1}^{a-1}q^j t^{a+b-j},&a>b.
  \end{cases}
\]
The second sum is empty when \(a=b+1\).  Thus a distinguished path below or
at the middle line contributes a positive symmetric interval.  A
distinguished path above the middle line contributes a negative open
interval.

\begin{conjecture}[Rational symmetric-string formula]
For every coprime \(r,s\ge2\),
\[
  \Cat_{r/s}(q,t)
  =
  \sum_{Q\in\T_{r/s}}
  \operatorname{Str}\bigl(\area(Q),\dinv(Q)\bigr).
\]
\end{conjecture}

The negative terms are not negative path counts.  They are cancellation terms
in a signed algebraic decomposition.  They occur precisely when a
distinguished path begins above the middle of its fixed-total-degree slice,
that is, when
\[
  \area(Q)>\dinv(Q)+1.
\]
They are common rather than exceptional: in the finite grid reported below,
80 of the 114 checked parameter pairs required nonzero negative corrections.

\section{What the checker does}

For each requested coprime pair \((r,s)\), \texttt{code.py} performs four
steps.

\begin{enumerate}
\item It generates every valid position-coordinate path.  The recursion is
      written in the weakly increasing path-height coordinates \(P_i\), so
      path validity is automatic.
\item It computes \(\area\), \(\defc\), and \(\dinv\), and forms the direct
      coefficient dictionary of \(\Cat_{r/s}(q,t)\).
\item It selects the paths with \(L_c=1\) and \(Q_c=0\), expands their signed
      strings, and forms the conjectural coefficient dictionary.
\item It compares the two dictionaries exactly.  It also checks the generated
      number of paths against the rational Catalan number
      \[
        \frac{1}{r+s}\binom{r+s}{r}.
      \]
\end{enumerate}

The implementation uses exact integer arithmetic and only the Python standard
library.  The deficit is accumulated while a path is generated, so shared
prefixes reuse earlier work rather than recomputing the full statistic for
every completed path.

\section{Finite verification}

The optimized Python checker was run on every ordered coprime pair satisfying
\[
  2\le r,s\le15.
\]
This consists of 114 cases and 11,597,208 generated rational Dyck paths.  Every
case passed the exact coefficient comparison.  Nonzero negative correction
terms occurred in 80 cases.  The two largest cases were
\[
  (r,s)=(14,15)\quad\text{and}\quad(15,14),
\]
each containing 2,674,440 paths.

Three small reference cases, written in the standard \((r,s)\) convention,
are
\[
\begin{array}{c|r|r|r}
(r,s)&\text{paths}&\text{positive terms}&\text{negative terms}\\
\hline
(5,3)&7&7&0\\
(8,5)&99&99&0\\
(12,7)&2652&2666&14
\end{array}
\]
The last case is the source paper's example written there in width/height
order as \((7,12)\).  On the development machine, the full grid completed in
about 25 seconds with two worker processes; runtime will vary by machine.

Useful commands are
\begin{verbatim}
python code.py
python code.py --case 5/3 --case 8/5 --case 12/7
python code.py --grid 2 15 --workers 2
\end{verbatim}
The default command checks \((r,s)=(12,7)\).

\section{Interpretation}

The verified grid gives substantial finite evidence, including many cases in
which cancellation by negative strings is genuinely required.  It does not
supply a uniform construction or proof for arbitrary coprime \((r,s)\).  The
main open problem recorded by this item is still to explain the signed-string
identity combinatorially in full generality.

\begin{thebibliography}{9}
\bibitem{Hawkes2024}
Graham Hawkes,
\emph{A Conjectured Formula for the Rational \(q,t\)-Catalan Polynomial},
Annals of Combinatorics 28 (2024), 749--795;
arXiv:2208.00577.
\end{thebibliography}

\end{document}
